\documentclass{article}

\usepackage[preprint]{neurips_2024}

\usepackage[utf8]{inputenc}
\usepackage[T1]{fontenc}
\usepackage{hyperref}
\usepackage{url}
\usepackage{booktabs}
\usepackage{amsmath,amssymb,amsfonts}
\usepackage{graphicx}
\usepackage{subcaption}
\usepackage{microtype}
\usepackage{nicefrac}
\usepackage{xcolor}
\usepackage{array}
\usepackage{makecell}
\usepackage{natbib}

\newcommand{\R}{\mathbb{R}}
\newcommand{\E}{\mathbb{E}}
\newcommand{\KL}{\mathrm{KL}}
\newcommand{\calN}{\mathcal{N}}
\newcommand{\bmu}{\boldsymbol{\mu}}
\newcommand{\bSigma}{\boldsymbol{\Sigma}}
\newcommand{\bpi}{\boldsymbol{\pi}}
\newcommand{\bx}{\mathbf{x}}
\newcommand{\by}{\mathbf{y}}
\newcommand{\bu}{\mathbf{u}}

\title{Reward-Conditioned Diffusion Scenarios for Portfolio Optimization}

\author{%
  Siddhant Sukhani \\
  Department of Management Science \& Engineering \\
  Stanford University \\
  \texttt{siddhantsukhani5@gmail.com} \\
  MS\&E 342: Stochastic Control and Optimization in Continuous Time
}

\begin{document}

\maketitle

\begin{abstract}
This report studies whether score-based diffusion scenarios can improve portfolio
construction after distributional calibration and reward-conditioned
fine-tuning. The pipeline trains a variance-preserving stochastic differential
equation score model on ten S\&P 500 sector ETF daily log returns from 2014 to
2020, evaluates hyperparameters only on 2021 validation returns, and reserves
2022 to 2024 for final testing. Gaussian Bures-Wasserstein optimal transport is
used to match generated first and second moments to training moments. A
KL-penalized reverse-diffusion control network is then trained with a
decision-aligned validation reward. The production evaluation passes the
leakage audit with zero failures and zero warnings. The strongest
fixed-scenario diffusion baseline obtains Sharpe 0.454,
while the validation-selected fine-tuned models underperform out of sample in
fixed mode. In rolling mode, all diffusion variants collapse toward rolling
Markowitz performance, with Sharpe ratios around 0.51. The results show that the
method is reproducible and mathematically coherent, but the current
fine-tuning reward overfits validation-period structure and does not yet provide
robust test-period alpha.
\end{abstract}

\section{Introduction}

Portfolio optimization depends on an estimate of the return distribution. The
classical Markowitz problem estimates means and covariances from historical
returns, then solves a mean-variance allocation problem \citep{markowitz1952}.
This is simple and transparent, but the estimated mean vector is noisy and the
resulting portfolio can become highly concentrated. Distributionally robust
portfolio optimization partially addresses this problem by optimizing against a
neighborhood of plausible distributions rather than a single empirical estimate
\citep{blanchet2022}.

Generative models offer a different route. Instead of estimating only moments,
they can produce synthetic return scenarios and let the downstream optimizer
solve a sample-average problem. Score-based diffusion models are a natural
candidate because they define a continuous-time noising process and learn the
reverse score needed for generation \citep{song2021}. Denoising diffusion models
have been successful in high-dimensional generative modeling \citep{ho2020}, and
the SDE formulation gives a direct link to stochastic control.

Financial returns, however, are not just generic multivariate data. They exhibit
stylized facts such as heavy tails, weak linear autocorrelation, volatility
clustering, and unstable dependence \citep{cont2001}. A return generator can look
reasonable under some marginal plots while still producing moments that imply
poor portfolios. The central question in this project is therefore practical:
can a diffusion model be made useful for portfolio decisions while preserving
enough distributional realism to avoid pure reward hacking?

The project evaluates three connected ideas. First, train a base
variance-preserving SDE model on historical sector ETF returns. Second, calibrate
generated scenarios with Gaussian optimal transport so that first and second
moments match the training distribution. Third, fine-tune the reverse diffusion
drift with a KL-penalized portfolio reward evaluated only on validation data. The
KL penalty is meant to limit how far the controlled reverse process can move from
the base model.

The final results are materially different from the preliminary positive result.
The production experiment evaluates both a base score model and an OT-augmented
score model, and selects eta separately for each family. Validation Sharpe is
very high, above 3.1 for the best candidates, but the selected fine-tuned
fixed-scenario strategies lose money during the test period. That mismatch is
the main empirical finding.

\section{Data and Experimental Protocol}

\subsection{Assets and Splits}

The dataset contains daily log returns for ten S\&P 500 sector ETFs: XLK, XLF,
XLE, XLV, XLI, XLY, XLP, XLU, XLB, and XLRE. The assets cover technology,
financials, energy, health care, industrials, consumer discretionary, consumer
staples, utilities, materials, and real estate. This universe is small enough for
controlled experimentation but broad enough to produce nontrivial covariance and
allocation structure.

The experiment uses a strict chronological split:
\begin{center}
\begin{tabular}{llll}
\toprule
Split & Date range & Trading days & Use \\
\midrule
Train & 2014-01-01 to 2020-12-31 & 1,317 & score model, scaler, OT target \\
Validation & 2021-01-01 to 2021-12-31 & 252 & eta selection, reward evaluation \\
Test & 2022-01-01 to 2024-12-31 & 756 & final performance evaluation \\
\bottomrule
\end{tabular}
\end{center}

The standardization scaler is fit only on training data. The experiment records
the training range, validation range, test-use status, and scaler provenance for
each model artifact. The final leakage audit reports a passing status, with zero
failures and zero warnings.

The chronological split is important because the test window contains a market
environment that differs sharply from the training window. The training period
includes the long post-crisis expansion, the 2018 volatility episode, and the
COVID-19 shock. The validation year, 2021, is a strong recovery year for many
risk assets. The test period includes the 2022 inflation and rate-hiking regime,
the 2023 rebound, and the 2024 continuation of concentrated equity strength.
This shift creates a difficult but useful evaluation setting. A model that
optimizes too directly for the validation year can select a portfolio that looks
excellent in 2021 but fails in the higher-rate regime that follows.

The sector ETF universe also makes the estimation problem concrete. The assets
are liquid, diversified within sector, and highly correlated during stress
periods. Their return distributions are not independent, and the covariance
matrix has economically meaningful structure. Technology and consumer
discretionary often carry high growth exposure; utilities and consumer staples
often provide more defensive exposure; energy has distinct commodity sensitivity.
The optimizer must therefore decide not only between high and low volatility
assets, but also between correlated factor exposures.

\subsection{Backtest Design}

The project compares equal-weight, rolling Markowitz, Wasserstein-robust
Markowitz, base diffusion scenarios, Gaussian-OT diffusion scenarios, and
fine-tuned diffusion scenarios. All reported test metrics use 2022 to 2024 data
only. The final comparison stage loads the frozen eta selection before test
evaluation and does not use test data for model selection.

Two backtest modes are reported. In fixed-scenario mode, each diffusion strategy
uses one scenario set throughout the test period. This isolates the scenario
generator as a static model of the training distribution. In rolling mode, each
monthly rebalance recalibrates scenarios to the trailing 252-day mean and
covariance using Gaussian OT. That mode tests whether diffusion scenarios add
anything after current rolling moments are imposed.

The baseline set is designed to separate three effects. Equal weight measures the
benefit of simple diversification without estimation. Rolling Markowitz measures
the value and instability of plug-in moment estimation. Wasserstein-robust
Markowitz measures whether shrinking toward a more conservative covariance
estimate improves realized risk. The diffusion strategies measure whether
scenario generation, moment calibration, and reward fine-tuning add incremental
value beyond those classical baselines. Because all strategies are evaluated on
the same daily return path, the comparison emphasizes allocation decisions rather
than differences in data handling.

\section{Diffusion Model}

\subsection{Forward SDE}

Let $X_0 \in \R^d$ denote a vector of standardized daily returns, with $d=10$.
The forward variance-preserving SDE is
\begin{equation}
  dX_t = -\frac{1}{2}\beta(t)X_t\,dt + \sqrt{\beta(t)}\,dW_t,
  \quad
  \beta(t)=\beta_{\min}+t(\beta_{\max}-\beta_{\min}),
\end{equation}
with $\beta_{\min}=0.1$ and $\beta_{\max}=20.0$. The conditional marginal is
\begin{equation}
  X_t|X_0 \sim \calN(\alpha_t X_0,\sigma_t^2 I),
  \quad
  \alpha_t=\exp\left(-\frac{1}{2}\int_0^t\beta(s)\,ds\right),
  \quad
  \sigma_t^2=1-\alpha_t^2.
\end{equation}

The score network $s_\theta(x,t)$ estimates
$\nabla_x \log p_t(x)$. Training minimizes denoising score matching:
\begin{equation}
  \mathcal{L}_{DSM}(\theta)
  =
  \E_{t,X_0,\epsilon}
  \left[
    \left\|
      s_\theta(\alpha_t X_0+\sigma_t\epsilon,t)
      + \epsilon/\sigma_t
    \right\|_2^2
  \right],
  \quad
  \epsilon \sim \calN(0,I).
\end{equation}

The network is a multilayer perceptron with hidden size 256, SiLU activations,
and a sinusoidal time embedding of dimension 64. The production model trains for
2,000 epochs with Adam at learning rate $10^{-3}$ \citep{kingma2015}. The final
training loss is 0.1574.

\subsection{Reverse Sampling}

Scenarios are generated by Euler-Maruyama discretization of the reverse process.
The production run uses 500 reverse steps and generates 10,000 base scenarios.
The generated samples are transformed back to return units using the
training-fitted scaler. The base scenario diagnostics record a
Bures-Wasserstein distance of 0.0309 from the training distribution and an
implied long-only MVO HHI near 1.0, meaning the raw generated moments induce an
almost fully concentrated portfolio.

This concentration problem is central. A generator can have acceptable
correlation plots while still producing a slightly biased mean vector. In
mean-variance optimization, small mean errors can dominate the allocation because
the inverse covariance amplifies the direction of perceived excess return. The
raw diffusion model underestimates volatility across assets and produces a
portfolio that is effectively a single-asset bet. That is not only undesirable
from a risk-management perspective; it is also evidence that distributional
diagnostics must be tied to downstream decisions. A scenario model should be
judged both by statistical realism and by the stability of the portfolio it
implies.

\subsection{Diagnostics Beyond Loss}

Training loss alone is not sufficient for this application. The project evaluates
generated scenarios along three dimensions. The first is marginal realism:
per-asset means, standard deviations, kurtosis, and quantile plots should be
close to the training returns. The second is cross-sectional realism: generated
correlations and covariances should preserve the dependence structure among
sectors. The third is decision realism: portfolios implied by generated
scenarios should not be pathologically concentrated or unstable.

These diagnostics reveal a clear pattern. The score model learns a reasonable
cross-sectional shape but compresses marginal volatility and fat tails. This is
common when a simple model is trained on limited financial data with a squared
denoising objective. The loss rewards average denoising accuracy, while the
tails are rare and carry less direct weight. For portfolio construction, this
matters because underestimated volatility can make the optimizer too aggressive.
Calibration is therefore not a cosmetic postprocessing step; it is required to
make the generated scenarios usable in a mean-variance allocation.

\section{Optimal Transport Calibration}

Gaussian OT is used as a moment-correcting postprocessor. Let generated samples
have mean and covariance $(\bmu_g,\bSigma_g)$, and let training returns have
$(\bmu_t,\bSigma_t)$. Under a Gaussian approximation, the Bures-Wasserstein map
from generated to target moments is affine:
\begin{equation}
  T(x)=\bmu_t + A(x-\bmu_g),
\end{equation}
where
\begin{equation}
  A = \bSigma_g^{-1/2}
      \left(
        \bSigma_g^{1/2}\bSigma_t\bSigma_g^{1/2}
      \right)^{1/2}
      \bSigma_g^{-1/2}.
\end{equation}
This construction follows the closed-form geometry of Gaussian optimal transport
and the Bures-Wasserstein metric \citep{bures1969,bhatia2019,peyre2019}. It
matches the first two moments exactly up to numerical error. It does not correct
tail behavior, skewness, volatility clustering, or time dependence.

This limitation is not a weakness of the implementation. It is a property of the
chosen map. A Gaussian OT map treats the source and target distributions as if
their relevant structure were fully summarized by mean and covariance. For a
mean-variance optimizer, that is a reasonable first-order correction because the
optimizer itself uses only those two moments. For a tail-risk optimizer, it would
be insufficient. In that setting, a calibration method would need to preserve or
reshape distributional features such as left-tail dependence and asymmetric
drawdowns.

The benefit of Gaussian OT is interpretability. If the calibrated scenarios
perform differently from the raw scenarios, the source of the difference is
clear: the first two moments have changed. This makes the fixed-scenario
comparison easier to diagnose. Raw diffusion performance reflects the generator's
uncontrolled moment bias. Gaussian-OT performance reflects a corrected
mean-covariance input to the same portfolio optimizer. Fine-tuned performance
then reflects a reward-conditioned change in scenario distribution.

The diagnostics show that Gaussian OT reduces the training-distance diagnostic
from 0.0309 to 0.00034 and reduces covariance Frobenius error from
0.00141 to $1.75\times10^{-5}$. The raw generated marginal standard deviations
are too low for every asset. For example, XLK has generated standard deviation
0.00558 versus training standard deviation 0.01486. After Gaussian OT, the XLK
standard deviation becomes 0.01476.

\begin{table}[t]
\centering
\caption{Scenario calibration diagnostics from the production evaluation.}
\label{tab:otdiag}
\begin{tabular}{lccc}
\toprule
Scenario set & W2 diagnostic & Covariance error & Implied HHI \\
\midrule
Base generated & 0.0309 & 0.00141 & 1.000 \\
Gaussian OT & 0.00034 & 0.000017 & 0.264 \\
Sinkhorn subset & diagnostic only & diagnostic only & not used \\
\bottomrule
\end{tabular}
\end{table}

\begin{figure}[t]
  \centering
  \includegraphics[width=0.92\textwidth]{figures/ot_gaussian_diagnostics.png}
  \caption{Gaussian OT diagnostics. The calibration aligns first and second
  moments to the training distribution, but it should not be interpreted as a
  full correction of non-Gaussian financial structure.}
  \label{fig:otdiag}
\end{figure}

\section{KL-Penalized Fine-Tuning}

\subsection{Controlled Reverse Process}

Fine-tuning adds a learnable drift perturbation $\bu_\phi(y,t)$ to the reverse
SDE:
\begin{equation}
  dY_t =
  \left[
    -\frac{1}{2}\beta(t)Y_t
    + \beta(t)s_\theta(Y_t,t)
    + \bu_\phi(Y_t,t)
  \right]dt
  + \sqrt{\beta(t)}\,dW_t.
\end{equation}
The control network is a two-hidden-layer MLP with hidden size 128. Its output is
bounded as $\tanh(\mathrm{MLP}(y,t))\cdot1.5$ to avoid extreme drifts.

By Girsanov's theorem, the path-space KL divergence between the controlled and
base reverse processes has a quadratic control-cost form \citep{oksendal2003}:
\begin{equation}
  \KL(P_\phi\|P_\theta)
  =
  \E_{P_\phi}
  \left[
    \int_0^T
    \frac{\|\bu_\phi(Y_t,t)\|_2^2}{2\beta(t)}
    dt
  \right].
\end{equation}
The fine-tuning objective is
\begin{equation}
  \max_{\phi}
  \E_{Y_0\sim P_\phi}[R(Y_0)]
  -
  \eta \KL(P_\phi\|P_\theta).
\end{equation}
Small $\eta$ permits larger drift changes. Large $\eta$ keeps the model closer
to the base generator.

\subsection{Portfolio Validation Reward}

For generated scenario batch $Y_0$, the method computes batch moments
$(\hat{\bmu},\hat{\bSigma})$ and a differentiable soft-Markowitz portfolio:
\begin{equation}
  \hat{\bpi}
  =
  \mathrm{softmax}\left((\hat{\bSigma}+10^{-4}I)^{-1}\hat{\bmu}\right).
\end{equation}
The softmax gives positive weights summing to one, but it is not the exact
solution to the constrained Markowitz problem. Exact differentiable convex
optimization layers would be a better future replacement \citep{agrawal2019}.

The softmax approximation has two practical consequences. First, it guarantees
long-only weights, which avoids short-sale leverage and keeps the allocation
comparable to the baseline constraints. Second, it can exaggerate small
differences in the unconstrained Markowitz direction. If one asset has a
slightly larger estimated risk-adjusted return, the softmax can push most of the
portfolio into that asset, especially when the covariance estimate is ill
conditioned. The concentration penalty is included to counteract this effect,
but the fixed-test results show that the penalty is not strong enough for the
most aggressive selected model.

The reward applies $\hat{\bpi}$ to fixed 2021 validation returns and combines
annualized mean, annualized variance, empirical CVaR, and concentration:
\begin{equation}
  R(Y_0)
  =
  252\bar{r}_v
  -
  \frac{\lambda}{2}252\,\mathrm{Var}(r_v)
  -
  \gamma \mathrm{CVaR}_{95}(r_v)
  -
  \tau \sum_i \hat{\pi}_i^2,
\end{equation}
where $\lambda=5$, $\gamma=0.5$, and $\tau=0.1$. Gradients flow through
$\hat{\bpi}$ to the generated scenario moments and back through the controlled
reverse process. Test returns are never used in this stage.

The reward is deliberately decision-aligned: generated scenarios are useful only
if the portfolio they imply performs well on validation returns. This is more
direct than matching stylized facts alone. It is also riskier. Validation returns
are a single historical realization, not a population objective. If 2021 rewards
a concentrated growth-oriented exposure, the fine-tuning step can learn to
generate scenarios that justify that exposure, even if the resulting portfolio
does not survive the next regime. The KL penalty limits the size of the drift
perturbation, but it does not eliminate validation overfitting.

\section{Eta Selection}

The eta sweep evaluates $\eta\in\{0.05,0.1,0.5,1.0,5.0\}$ for both the base score
model and the OT-augmented score model. Selection is based only on 2021
validation Sharpe. The selection protocol explicitly records that test metrics
are not available at selection time.

\begin{table}[t]
\centering
\caption{Validation eta sweep. The selected eta is the highest validation Sharpe
within each model family.}
\label{tab:eta}
\begin{tabular}{llcccc}
\toprule
Model & Eta & Val. Sharpe & Ann. return & Ann. vol. & HHI \\
\midrule
base & 0.05 & 3.136 & 46.04\% & 14.68\% & 0.982 \\
base & 0.10 & 3.128 & 46.10\% & 14.74\% & 1.000 \\
base & 0.50 & 2.921 & 41.49\% & 14.20\% & 0.447 \\
base & 1.00 & 2.086 & 36.49\% & 17.49\% & 0.732 \\
base & 5.00 & 1.805 & 34.74\% & 19.25\% & 1.000 \\
\midrule
OT-aug. & 0.05 & 3.224 & 44.08\% & 13.67\% & 0.714 \\
OT-aug. & 0.10 & 3.139 & 48.74\% & 15.53\% & 0.532 \\
OT-aug. & 0.50 & 2.042 & 34.97\% & 17.12\% & 0.684 \\
OT-aug. & 1.00 & \textbf{3.234} & 41.65\% & 12.88\% & 0.520 \\
OT-aug. & 5.00 & 2.071 & 34.16\% & 16.49\% & 0.587 \\
\bottomrule
\end{tabular}
\end{table}

The selected models are OT-augmented eta 1.0 and base eta 0.05. Both selection
scores are high, but both turn out to be poor fixed-scenario test portfolios.
This gap is evidence that the reward and eta selection protocol can overfit a
single validation year.

The eta table also shows that validation Sharpe alone is not a sufficient model
selection criterion. The base eta 0.05 and eta 0.10 models both reach validation
Sharpe above 3.1, but their HHI values are essentially one. A selection rule that
jointly constrained validation Sharpe, HHI, and sensitivity to small changes in
the validation window would likely reject those models. The OT-augmented eta 1.0
model is less concentrated, but its fixed-test return is still negative. This
suggests that concentration is not the only failure mode. The validation reward
can also select a sector tilt that is diversified but regime-specific.

\begin{figure}[t]
  \centering
  \includegraphics[width=0.78\textwidth]{figures/eta_validation_frontier.png}
  \caption{Validation frontier used for eta selection. Test data are not used in
  this figure or in the eta-selection procedure.}
  \label{fig:eta}
\end{figure}

\section{Fixed-Scenario Results}

Table \ref{tab:fixed} reports the fixed-scenario test results. Equal weight earns
5.59\% annually with Sharpe 0.363. Rolling Markowitz earns 9.03\% with Sharpe
0.515, but it is concentrated with HHI 0.541 and high turnover. The
Wasserstein-robust baseline is diversified and low turnover, but its Sharpe is
only 0.380.

Among diffusion strategies, the raw base scenarios produce annual return 11.48\%
and Sharpe 0.454, but with annual volatility 25.31\%, maximum drawdown 33.13\%,
and HHI 1.0. This is not a usable allocation profile. Gaussian OT lowers
volatility and concentration, but also lowers return, producing Sharpe 0.393.
The validation-selected fine-tuned models perform worst in fixed mode:
OT-augmented eta 1.0 has Sharpe -0.059, and base eta 0.05 has Sharpe -0.221.

\begin{table}[t]
\centering
\caption{Fixed-scenario backtest, 2022 to 2024. Sharpe confidence intervals use
2,000 i.i.d. bootstrap draws.}
\label{tab:fixed}
\setlength{\tabcolsep}{4pt}
\begin{tabular}{lcccccc}
\toprule
Strategy & Ann. ret. & Ann. vol. & Sharpe & 95\% CI & Max DD & HHI \\
\midrule
Equal-Weight & 5.59\% & 15.41\% & 0.363 & $[-0.76,1.46]$ & -18.41\% & 0.100 \\
Markowitz & 9.03\% & 17.54\% & 0.515 & $[-0.60,1.57]$ & -17.09\% & 0.541 \\
Wasserstein-Robust & 5.75\% & 15.12\% & 0.380 & $[-0.74,1.48]$ & -17.92\% & 0.103 \\
Diffusion-Base & 11.48\% & 25.31\% & 0.454 & $[-0.66,1.55]$ & -33.13\% & 1.000 \\
Diffusion+GaussOT & 6.50\% & 16.56\% & 0.393 & $[-0.69,1.50]$ & -22.97\% & 0.264 \\
Diffusion+FT, OT-aug. eta 1.0 & -1.08\% & 18.35\% & -0.059 & $[-1.18,1.06]$ & -28.82\% & 0.520 \\
Diffusion+FT, base eta 0.05 & -4.63\% & 20.93\% & -0.221 & $[-1.37,0.88]$ & -33.20\% & 0.982 \\
\bottomrule
\end{tabular}
\end{table}

\begin{figure}[t]
  \centering
  \includegraphics[width=\textwidth]{figures/final_comparison_fixed.png}
  \caption{Fixed-scenario test comparison. The fine-tuned strategies selected
  by validation Sharpe do not generalize in fixed-scenario mode.}
  \label{fig:fixed}
\end{figure}

The confidence intervals are wide. That is expected with about 756 daily test
observations, especially when returns are non-Gaussian and serial dependence is
possible. The simple bootstrap is useful as a warning against overclaiming, but a
stationary or block bootstrap would be more appropriate for dependent financial
returns \citep{politis1994}. No pairwise Sharpe difference should be treated as
statistically decisive.

The fixed-scenario results are still informative. They show that the base
diffusion model contains enough signal to generate a high-return allocation, but
not enough calibration to control risk. They also show that first-two-moment
calibration improves risk characteristics but does not create superior realized
performance. Most importantly, they show that validation reward optimization can
move the model in the wrong direction. This is the central empirical warning:
decision-aligned fine-tuning is powerful enough to change the portfolio, but the
current validation design is not strong enough to guarantee that the change is
robust.

\section{Rolling Results}

Table \ref{tab:rolling} reports the rolling-diffusion test results. Once
generated scenarios are recalibrated to the trailing 252-day historical moments,
all diffusion variants produce nearly the same performance as rolling Markowitz.
Diffusion-Base obtains Sharpe 0.5105, Diffusion+GaussOT obtains 0.5096, the
OT-augmented fine-tuned model obtains 0.5101, and the base fine-tuned model
obtains 0.5106. These values are economically indistinguishable from rolling
Markowitz at 0.5151.

\begin{table}[t]
\centering
\caption{Rolling-diffusion backtest, 2022 to 2024. Scenarios are recalibrated
to trailing 252-day moments at each rebalance using only past data.}
\label{tab:rolling}
\begin{tabular}{lccccc}
\toprule
Strategy & Ann. ret. & Ann. vol. & Sharpe & Max DD & Turnover \\
\midrule
Equal-Weight & 5.59\% & 15.41\% & 0.363 & -18.41\% & 0.000 \\
Markowitz & 9.03\% & 17.54\% & 0.515 & -17.09\% & 0.617 \\
Wasserstein-Robust & 5.75\% & 15.12\% & 0.380 & -17.92\% & 0.040 \\
Diffusion-Base & 8.90\% & 17.44\% & 0.511 & -17.10\% & 0.588 \\
Diffusion+GaussOT & 8.87\% & 17.40\% & 0.510 & -17.09\% & 0.586 \\
Diffusion+FT, OT-aug. eta 1.0 & 8.88\% & 17.41\% & 0.510 & -17.09\% & 0.585 \\
Diffusion+FT, base eta 0.05 & 8.90\% & 17.43\% & 0.511 & -17.10\% & 0.588 \\
\bottomrule
\end{tabular}
\end{table}

\begin{figure}[t]
  \centering
  \includegraphics[width=\textwidth]{figures/final_comparison_rolling.png}
  \caption{Rolling-diffusion test comparison. Moment recalibration makes
  diffusion strategies behave almost like rolling Markowitz.}
  \label{fig:rolling}
\end{figure}

This collapse has a clear explanation. Mean-variance optimization depends only
on the estimated mean and covariance. Rolling Gaussian OT forces the scenario set
to match the same trailing moments that rolling Markowitz already uses. Unless
the portfolio objective uses higher-order scenario information, or the generator
contributes predictive structure beyond trailing moments, the rolling diffusion
strategy has little room to differ.

This result should not be read as a failure of diffusion modeling in general. It
is a consequence of pairing a moment-matched scenario generator with a
mean-variance optimizer. Once the rolling calibration step imposes the current
mean and covariance, the remaining higher-order differences in the scenario set
are mostly invisible to the optimizer. A different downstream objective could
use those differences. For example, a CVaR optimizer would depend on the lower
tail of the generated distribution. A drawdown-aware objective would depend on
path structure. A regime-conditioned optimizer could use the generator to
forecast state-dependent covariance rather than simply reproduce trailing
moments.

\section{Interpretation}

\subsection{What Worked}

The pipeline is reproducible and properly separated by time. The audit passed.
The score model generated finite scenarios with recorded provenance. Gaussian OT
did exactly what it was supposed to do: it aligned first and second moments and
reduced concentration relative to raw generated scenarios. The rolling test also
behaved consistently with the mathematics of the pipeline.

The project also demonstrates a useful engineering structure. Every stage is
automated, including data preparation, leakage audit, score training, sampling,
OT calibration, eta sweep, comparison, and report asset generation. This makes
the experiment repeatable without changing the modeling protocol between runs.

\subsection{What Did Not Work}

The fine-tuning reward did not generalize. Validation Sharpe selected highly
aggressive portfolios. The base eta 0.05 model has validation Sharpe 3.136 and
HHI 0.982, then loses 4.63\% annually in fixed test with HHI 0.982. That is a
classic sign of validation-period overfitting and excessive concentration. The
OT-augmented eta 1.0 model is less concentrated, but it still has negative
fixed-test return.

The base generator also underestimates marginal volatility and fat tails before
OT calibration. The generated kurtosis values in the diagnostics are near zero,
while training-return kurtosis is high across assets. Gaussian OT fixes
volatility but cannot fix this tail mismatch. As a result, the scenario
distribution is moment-aligned but not fully realistic.

\subsection{Why the Earlier Positive Claim Changed}

The prior draft stated that fine-tuning achieved Sharpe 0.70 in fixed mode. The
final metrics do not support that statement. The updated selection protocol
evaluates both base and OT-augmented model families and chooses different etas.
Under those frozen selections, fixed-test fine-tuning underperforms. The report
therefore treats the final production evaluation as authoritative.

\section{Limitations}

The first limitation is model expressiveness. A plain MLP score network on
one-day return vectors does not model temporal dynamics directly. It sees the
cross-section of daily returns, not a sequence of market states. That makes it
difficult to capture volatility clustering and regime dependence, both of which
matter in financial data.

The second limitation is the reward. The soft-Markowitz reward is differentiable,
but it can produce highly concentrated portfolios when the generated moment
estimates are noisy. It optimizes validation-year performance rather than a
more stable distributional objective. Stronger concentration controls,
turnover-aware rewards, and nested cross-validation across multiple validation
regimes would make eta selection more reliable.

The third limitation is the calibration layer. Gaussian OT is useful and
transparent, but it reduces scenario calibration to moment matching. If the
downstream optimizer is mean-variance, then moment matching may be enough. If
the project goal is to exploit higher-order scenario structure, the optimizer
must use tail or path-dependent objectives, and the generator must be evaluated
on those properties.

The fourth limitation is statistical power. A three-year test window is small
for Sharpe-ratio inference. The confidence intervals in Tables \ref{tab:fixed}
and \ref{tab:rolling} are wide even under a simple i.i.d. bootstrap. A longer
history, walk-forward validation, and block bootstrap inference are needed before
making strong performance claims.

The fifth limitation is the objective mismatch between training diagnostics and
final use. The score model is trained to estimate a denoising score, the OT map
is evaluated by moment alignment, the fine-tuning step is evaluated by
validation Sharpe, and the final report evaluates realized test performance.
Each stage is reasonable in isolation, but the objectives are not identical.
Future work should reduce this mismatch by defining a single decision-aware
evaluation protocol that is used consistently for model selection and final
assessment.

\section{Future Work}

The most important next step is to replace single-year eta selection with a
walk-forward validation design. For example, train on 2014 to 2018, validate on
2019, train on 2014 to 2019, validate on 2020, then reserve 2021 for final
model selection before testing on 2022 to 2024. This would reduce dependence on
one unusually favorable validation year.

The second improvement is to use a differentiable convex optimization layer for
the portfolio decision. That would align the training reward with the exact
long-only constrained optimizer used in evaluation, rather than relying on a
softmax approximation.

The third improvement is to condition the diffusion model on market state. A
regime-conditioned score model could generate scenarios conditional on realized
volatility, trend, rates, or macro indicators. That would give the rolling
strategy a way to differ from simple trailing-moment Markowitz.

The fourth improvement is to evaluate objectives that actually use scenario
tails. If the downstream portfolio remains mean-variance, the generator is judged
mostly through means and covariances. CVaR optimization, drawdown-aware
objectives, or utility functions with downside asymmetry would better test
whether diffusion scenarios add value beyond moment estimation.

\section{Conclusion}

The latest production run supports a sober conclusion. The project successfully
builds a controlled, reproducible pipeline for diffusion-based portfolio
scenarios with strict data separation, Gaussian OT calibration, and
KL-penalized reward fine-tuning. The mathematical method is sound and the
calibration diagnostics behave as expected. The leakage audit passes cleanly.

The empirical result is mixed. Base diffusion scenarios have high fixed-test
return but unacceptable concentration and drawdown. Gaussian OT improves moment
alignment and concentration, but it does not beat rolling Markowitz. The
validation-selected fine-tuned models do not generalize in fixed-scenario mode.
In rolling mode, Gaussian moment recalibration causes diffusion strategies to
converge to rolling Markowitz performance.

The main lesson is that reward-conditioned diffusion for portfolios is promising
as a framework, but the current implementation is not yet a robust trading
method. The next version should focus on stronger validation, exact
differentiable portfolio optimization, regime conditioning, and tail-aware
objectives.

\clearpage
\section*{References}

\begin{thebibliography}{99}

\bibitem[Agrawal et al.(2019)]{agrawal2019}
Agrawal, A., Amos, B., Barratt, S., Boyd, S., Diamond, S., and Kolter, J. Z.
(2019).
\newblock Differentiable convex optimization layers.
\newblock \emph{Advances in Neural Information Processing Systems}, 32.
\newblock \url{https://papers.neurips.cc/paper/9152-differentiable-convex-optimization-layers}.

\bibitem[Anderson(1982)]{anderson1982}
Anderson, B. D. O. (1982).
\newblock Reverse-time diffusion equation models.
\newblock \emph{Stochastic Processes and Their Applications}, 12(3):313--326.

\bibitem[Bhatia et al.(2019)]{bhatia2019}
Bhatia, R., Jain, T., and Lim, Y. (2019).
\newblock On the Bures-Wasserstein distance between positive definite matrices.
\newblock \emph{Expositiones Mathematicae}, 37(2):165--191.
\newblock \url{https://arxiv.org/abs/1712.01504}.

\bibitem[Blanchet et al.(2022)]{blanchet2022}
Blanchet, J., Chen, L., and Zhou, X. Y. (2022).
\newblock Distributionally robust mean-variance portfolio selection with
Wasserstein distances.
\newblock \emph{Management Science}, 68(9):6382--6410.
\newblock \url{https://pubsonline.informs.org/doi/abs/10.1287/mnsc.2021.4155}.

\bibitem[Bures(1969)]{bures1969}
Bures, D. (1969).
\newblock An extension of Kakutani's theorem on infinite product measures to the
tensor product of semifinite $W^*$-algebras.
\newblock \emph{Transactions of the American Mathematical Society}, 135:199--212.

\bibitem[Cont(2001)]{cont2001}
Cont, R. (2001).
\newblock Empirical properties of asset returns: Stylized facts and statistical
issues.
\newblock \emph{Quantitative Finance}, 1(2):223--236.
\newblock \url{https://www.tandfonline.com/doi/abs/10.1080/713665670}.

\bibitem[Ho et al.(2020)]{ho2020}
Ho, J., Jain, A., and Abbeel, P. (2020).
\newblock Denoising diffusion probabilistic models.
\newblock \emph{Advances in Neural Information Processing Systems}, 33:6840--6851.
\newblock \url{https://proceedings.neurips.cc/paper/2020/hash/4c5bcfec8584af0d967f1ab10179ca4b-Abstract.html}.

\bibitem[Kingma and Ba(2015)]{kingma2015}
Kingma, D. P. and Ba, J. (2015).
\newblock Adam: A method for stochastic optimization.
\newblock \emph{International Conference on Learning Representations}.

\bibitem[Markowitz(1952)]{markowitz1952}
Markowitz, H. (1952).
\newblock Portfolio selection.
\newblock \emph{The Journal of Finance}, 7(1):77--91.

\bibitem[Oksendal(2003)]{oksendal2003}
Oksendal, B. (2003).
\newblock \emph{Stochastic Differential Equations: An Introduction with
Applications}.
\newblock Springer, 6th edition.

\bibitem[Peyre and Cuturi(2019)]{peyre2019}
Peyre, G. and Cuturi, M. (2019).
\newblock Computational optimal transport.
\newblock \emph{Foundations and Trends in Machine Learning}, 11(5--6):355--607.
\newblock \url{https://arxiv.org/abs/1803.00567}.

\bibitem[Politis and Romano(1994)]{politis1994}
Politis, D. N. and Romano, J. P. (1994).
\newblock The stationary bootstrap.
\newblock \emph{Journal of the American Statistical Association}, 89(428):1303--1313.
\newblock \url{https://www.tandfonline.com/doi/abs/10.1080/01621459.1994.10476870}.

\bibitem[Song et al.(2021)]{song2021}
Song, Y., Sohl-Dickstein, J., Kingma, D. P., Kumar, A., Ermon, S., and Poole, B.
(2021).
\newblock Score-based generative modeling through stochastic differential
equations.
\newblock \emph{International Conference on Learning Representations}.
\newblock \url{https://research.google/pubs/score-based-generative-modeling-through-stochastic-differential-equations/}.

\end{thebibliography}

\end{document}
